% =============================================================================
% appendix.tex  —  % =============================================================================
% appendix.tex  —  % =============================================================================
% appendix.tex  —  % =============================================================================
% appendix.tex  —  \input{appendix} before \end{document} in discoure.tex
% Provides the replication appendices the reviewers asked for:
%   A. 26-term thematic keyword list (W1 / reproducibility)
%   B. Annotation codebook summary (reproducibility)
%   C. LLM annotation prompt + few-shot setup (reproducibility)
% =============================================================================

\appendices

% -----------------------------------------------------------------------------
\section{Thematic Keyword List}
\label{app:keywords}
% -----------------------------------------------------------------------------
The thematic keyword pass (Section~III.B) matches the following 26 terms
case-insensitively against the full cleaned text. A document is retained if it
matches at least one term. The list is reproduced in full here so that the
inclusion/exclusion decision is auditable; the prose summary in the body groups
these terms into three categories.

\begin{table}[!ht]
\renewcommand{\arraystretch}{1.2}
\caption{The 26 earthquake-relevance keywords used at the corpus-construction
stage. Matching is case-insensitive substring matching on the cleaned text.}
\label{table_keywords}
\centering
\footnotesize
\begin{tabular}{|l|l|}
\hline
\textbf{Category} & \textbf{Terms} \\
\hline
Core disaster &
\begin{tabular}[t]{@{}l@{}}
deprem, enkaz, hasar, y\i k\i m,\\
depremzede, 6 \c{S}ubat,\\
kahramanmara\c{s} depremi
\end{tabular} \\
\hline
Institutional / operational &
\begin{tabular}[t]{@{}l@{}}
AFAD, kurtarma, tahliye,\\
acil durum, \c{c}ad\i r kent,\\
bar\i nak, konteyner, prefabrik
\end{tabular} \\
\hline
Reconstruction &
\begin{tabular}[t]{@{}l@{}}
kal\i c\i\ konut, ge\c{c}ici bar\i nma,\\
yeniden yap\i m, yeniden yap\i lanma,\\
d\"on\"u\c{s}\"um, kentsel d\"on\"u\c{s}\"um,\\
TOK\.{I}, in\c{s}aat, altyap\i
\end{tabular} \\
\hline
\end{tabular}
\end{table}

This keyword pass should be understood as a deliberate scoping rule rather than a
neutral recovery mechanism for all earthquake-relevant discourse. It produces a
conservative analytic corpus centred on explicit disaster and reconstruction
vocabulary, but it may systematically exclude texts that discuss the earthquake
indirectly, metaphorically, or through less canonical local terminology. The
downstream results are therefore claims about this explicitly earthquake-marked
official discourse, not about every possible official text whose meaning may be
connected to the disaster.

% -----------------------------------------------------------------------------
\section{Annotation Codebook (Summary)}
\label{app:codebook}
% -----------------------------------------------------------------------------
The full codebook is included in the replication repository
(\texttt{annotation/codebook.md}). The operational core is summarised below.
Every frame is scored independently on a 0--3 ordinal scale; a single document may
score high on more than one frame.

\noindent\textbf{Decision procedure (applied per frame).}
\begin{enumerate}
\item \emph{Are the frame's indicators present at all?} No $\rightarrow$ 0.
\item \emph{Is the frame central or peripheral?} Peripheral (1--2 sentences)
      $\rightarrow$ 1.
\item \emph{Is the frame the backbone of the text?} Covers $\sim$30--50\%
      $\rightarrow$ 2; title + lead + conclusion all carry the frame $\rightarrow$ 3.
\end{enumerate}

\noindent\textbf{Frame definitions and indicators.}
\begin{itemize}
\item \textbf{Technical resilience} --- framing reconstruction through
engineering, building codes, inspection, materials, soil surveys, and damage
assessment. Indicators: \textit{yap\i\ denetimi, zemin et\"ud\"u, y\"onetmelik,
kentsel d\"on\"u\c{s}\"um} (engineering sense), \textit{g\"u\c{c}lendirme, AFAD
raporu}, numerical/technical detail.
\item \textbf{Political mobilisation} --- framing reconstruction through electoral
gain, partisan competition, leadership legitimacy, and national-mobilisation
rhetoric. Indicators: \textit{milletimize s\"oz verdik, muhalefet engelledi, Say\i n
Cumhurba\c{s}kan\i m\i z, AK Parti}, campaign pledges, ceremonial register.
\item \textbf{Developmental transformation} --- framing the disaster as a
development opportunity, investment, infrastructure modernisation, or regional
growth. Indicators: \textit{yeni \c{s}ehir, modern altyap\i, yat\i r\i m, TOK\.{I}
konutlar\i, OSB, istihdam}, project emphasis.
\item \textbf{Sustainability} --- framing reconstruction through climate
resilience, green building, green infrastructure, or ecological-city concepts;
rare in this corpus. Indicators: \textit{ye\c{s}il bina, iklim diren\c{c}li,
ekolojik, karbon ayak izi, yenilenebilir enerji, enerji verimlili\u{g}i}. Note:
``modern, comfortable housing'' is \emph{not} sustainability absent an explicit
environmental term.
\end{itemize}

% -----------------------------------------------------------------------------
\section{LLM Annotation Prompt}
\label{app:llmprompt}
% -----------------------------------------------------------------------------
The 201 documents outside the human double-coded subset were annotated via the
Anthropic API at temperature 0, with \texttt{max\_tokens} bounded and the model
asked to return a strict JSON object with the four frame scores. The prompt
combines (i) the codebook block reproduced verbatim below, (ii) eight few-shot
examples drawn from the human-agreement subset, and (iii) the target document
(truncated to a fixed character budget). The exact model identifier and API
version used are recorded in the configuration cell of the annotation notebook
(\texttt{annotation/LLM\_annotation.ipynb}). Two Anthropic models appear in the
notebook by design: an initial Claude Haiku 4.5 trial used for model selection,
and Claude Sonnet 4.5, which produced all 201 deployment-subset labels in the
final dataset. Haiku contributed no labels to the training set
(see Section~IV and the model-selection paragraph). The configuration cell pins
the exact model strings (\texttt{claude-sonnet-4-5} and
\texttt{claude-haiku-4-5-20251001}); these should match the prose verbatim.}

\noindent\textbf{Codebook block (verbatim, Turkish).}
\begin{quote}\footnotesize\ttfamily
4 frame'i 0-3 Likert \"ol\c{c}e\u{g}inde de\u{g}erlendir. Her frame BA\u{G}IMSIZ
de\u{g}erlendirilir; bir metin birden fazla frame'den y\"uksek puan alabilir.

FRAME 1 --- TECHNICAL; FRAME 2 --- POLITICAL; FRAME 3 --- DEVELOPMENT;
FRAME 4 --- SUSTAINABILITY (definitions and indicators as in
Appendix~\ref{app:codebook}).

KARAR REHBER\.{I}:
1. Frame'in g\"ostergeleri VAR MI? Hay\i r $\rightarrow$ 0. Evet:
2. Metnin MERKEZ\.{I}NDE M\.{I}? K\i y\i\ (1-2 c\"umle) $\rightarrow$ 1. Merkez:
3. Metnin OMURGASI MI? Hay\i r ($\sim$\%30-50) $\rightarrow$ 2.
   Evet (ba\c{s}l\i k+giri\c{s}+sonu\c{c}) $\rightarrow$ 3.
\end{quote}

\noindent The model was instructed to output only a JSON object of the form
\texttt{\{"technical": <0-3>, "political": <0-3>, "development": <0-3>,
"sustainability": <0-3>\}}, which was parsed and range-checked
(0$\le v \le$3) before being written to the label store with a
\texttt{label\_source} of \texttt{llm\_only} or, after manual review,
\texttt{llm\_then\_manual}.
 before \end{document} in discoure.tex
% Provides the replication appendices the reviewers asked for:
%   A. 26-term thematic keyword list (W1 / reproducibility)
%   B. Annotation codebook summary (reproducibility)
%   C. LLM annotation prompt + few-shot setup (reproducibility)
% =============================================================================

\appendices

% -----------------------------------------------------------------------------
\section{Thematic Keyword List}
\label{app:keywords}
% -----------------------------------------------------------------------------
The thematic keyword pass (Section~III.B) matches the following 26 terms
case-insensitively against the full cleaned text. A document is retained if it
matches at least one term. The list is reproduced in full here so that the
inclusion/exclusion decision is auditable; the prose summary in the body groups
these terms into three categories.

\begin{table}[!ht]
\renewcommand{\arraystretch}{1.2}
\caption{The 26 earthquake-relevance keywords used at the corpus-construction
stage. Matching is case-insensitive substring matching on the cleaned text.}
\label{table_keywords}
\centering
\footnotesize
\begin{tabular}{|l|l|}
\hline
\textbf{Category} & \textbf{Terms} \\
\hline
Core disaster &
\begin{tabular}[t]{@{}l@{}}
deprem, enkaz, hasar, y\i k\i m,\\
depremzede, 6 \c{S}ubat,\\
kahramanmara\c{s} depremi
\end{tabular} \\
\hline
Institutional / operational &
\begin{tabular}[t]{@{}l@{}}
AFAD, kurtarma, tahliye,\\
acil durum, \c{c}ad\i r kent,\\
bar\i nak, konteyner, prefabrik
\end{tabular} \\
\hline
Reconstruction &
\begin{tabular}[t]{@{}l@{}}
kal\i c\i\ konut, ge\c{c}ici bar\i nma,\\
yeniden yap\i m, yeniden yap\i lanma,\\
d\"on\"u\c{s}\"um, kentsel d\"on\"u\c{s}\"um,\\
TOK\.{I}, in\c{s}aat, altyap\i
\end{tabular} \\
\hline
\end{tabular}
\end{table}

This keyword pass should be understood as a deliberate scoping rule rather than a
neutral recovery mechanism for all earthquake-relevant discourse. It produces a
conservative analytic corpus centred on explicit disaster and reconstruction
vocabulary, but it may systematically exclude texts that discuss the earthquake
indirectly, metaphorically, or through less canonical local terminology. The
downstream results are therefore claims about this explicitly earthquake-marked
official discourse, not about every possible official text whose meaning may be
connected to the disaster.

% -----------------------------------------------------------------------------
\section{Annotation Codebook (Summary)}
\label{app:codebook}
% -----------------------------------------------------------------------------
The full codebook is included in the replication repository
(\texttt{annotation/codebook.md}). The operational core is summarised below.
Every frame is scored independently on a 0--3 ordinal scale; a single document may
score high on more than one frame.

\noindent\textbf{Decision procedure (applied per frame).}
\begin{enumerate}
\item \emph{Are the frame's indicators present at all?} No $\rightarrow$ 0.
\item \emph{Is the frame central or peripheral?} Peripheral (1--2 sentences)
      $\rightarrow$ 1.
\item \emph{Is the frame the backbone of the text?} Covers $\sim$30--50\%
      $\rightarrow$ 2; title + lead + conclusion all carry the frame $\rightarrow$ 3.
\end{enumerate}

\noindent\textbf{Frame definitions and indicators.}
\begin{itemize}
\item \textbf{Technical resilience} --- framing reconstruction through
engineering, building codes, inspection, materials, soil surveys, and damage
assessment. Indicators: \textit{yap\i\ denetimi, zemin et\"ud\"u, y\"onetmelik,
kentsel d\"on\"u\c{s}\"um} (engineering sense), \textit{g\"u\c{c}lendirme, AFAD
raporu}, numerical/technical detail.
\item \textbf{Political mobilisation} --- framing reconstruction through electoral
gain, partisan competition, leadership legitimacy, and national-mobilisation
rhetoric. Indicators: \textit{milletimize s\"oz verdik, muhalefet engelledi, Say\i n
Cumhurba\c{s}kan\i m\i z, AK Parti}, campaign pledges, ceremonial register.
\item \textbf{Developmental transformation} --- framing the disaster as a
development opportunity, investment, infrastructure modernisation, or regional
growth. Indicators: \textit{yeni \c{s}ehir, modern altyap\i, yat\i r\i m, TOK\.{I}
konutlar\i, OSB, istihdam}, project emphasis.
\item \textbf{Sustainability} --- framing reconstruction through climate
resilience, green building, green infrastructure, or ecological-city concepts;
rare in this corpus. Indicators: \textit{ye\c{s}il bina, iklim diren\c{c}li,
ekolojik, karbon ayak izi, yenilenebilir enerji, enerji verimlili\u{g}i}. Note:
``modern, comfortable housing'' is \emph{not} sustainability absent an explicit
environmental term.
\end{itemize}

% -----------------------------------------------------------------------------
\section{LLM Annotation Prompt}
\label{app:llmprompt}
% -----------------------------------------------------------------------------
The 201 documents outside the human double-coded subset were annotated via the
Anthropic API at temperature 0, with \texttt{max\_tokens} bounded and the model
asked to return a strict JSON object with the four frame scores. The prompt
combines (i) the codebook block reproduced verbatim below, (ii) eight few-shot
examples drawn from the human-agreement subset, and (iii) the target document
(truncated to a fixed character budget). The exact model identifier and API
version used are recorded in the configuration cell of the annotation notebook
(\texttt{annotation/LLM\_annotation.ipynb}). Two Anthropic models appear in the
notebook by design: an initial Claude Haiku 4.5 trial used for model selection,
and Claude Sonnet 4.5, which produced all 201 deployment-subset labels in the
final dataset. Haiku contributed no labels to the training set
(see Section~IV and the model-selection paragraph). The configuration cell pins
the exact model strings (\texttt{claude-sonnet-4-5} and
\texttt{claude-haiku-4-5-20251001}); these should match the prose verbatim.}

\noindent\textbf{Codebook block (verbatim, Turkish).}
\begin{quote}\footnotesize\ttfamily
4 frame'i 0-3 Likert \"ol\c{c}e\u{g}inde de\u{g}erlendir. Her frame BA\u{G}IMSIZ
de\u{g}erlendirilir; bir metin birden fazla frame'den y\"uksek puan alabilir.

FRAME 1 --- TECHNICAL; FRAME 2 --- POLITICAL; FRAME 3 --- DEVELOPMENT;
FRAME 4 --- SUSTAINABILITY (definitions and indicators as in
Appendix~\ref{app:codebook}).

KARAR REHBER\.{I}:
1. Frame'in g\"ostergeleri VAR MI? Hay\i r $\rightarrow$ 0. Evet:
2. Metnin MERKEZ\.{I}NDE M\.{I}? K\i y\i\ (1-2 c\"umle) $\rightarrow$ 1. Merkez:
3. Metnin OMURGASI MI? Hay\i r ($\sim$\%30-50) $\rightarrow$ 2.
   Evet (ba\c{s}l\i k+giri\c{s}+sonu\c{c}) $\rightarrow$ 3.
\end{quote}

\noindent The model was instructed to output only a JSON object of the form
\texttt{\{"technical": <0-3>, "political": <0-3>, "development": <0-3>,
"sustainability": <0-3>\}}, which was parsed and range-checked
(0$\le v \le$3) before being written to the label store with a
\texttt{label\_source} of \texttt{llm\_only} or, after manual review,
\texttt{llm\_then\_manual}.
 before \end{document} in discoure.tex
% Provides the replication appendices the reviewers asked for:
%   A. 26-term thematic keyword list (W1 / reproducibility)
%   B. Annotation codebook summary (reproducibility)
%   C. LLM annotation prompt + few-shot setup (reproducibility)
% =============================================================================

\appendices

% -----------------------------------------------------------------------------
\section{Thematic Keyword List}
\label{app:keywords}
% -----------------------------------------------------------------------------
The thematic keyword pass (Section~III.B) matches the following 26 terms
case-insensitively against the full cleaned text. A document is retained if it
matches at least one term. The list is reproduced in full here so that the
inclusion/exclusion decision is auditable; the prose summary in the body groups
these terms into three categories.

\begin{table}[!ht]
\renewcommand{\arraystretch}{1.2}
\caption{The 26 earthquake-relevance keywords used at the corpus-construction
stage. Matching is case-insensitive substring matching on the cleaned text.}
\label{table_keywords}
\centering
\footnotesize
\begin{tabular}{|l|l|}
\hline
\textbf{Category} & \textbf{Terms} \\
\hline
Core disaster &
\begin{tabular}[t]{@{}l@{}}
deprem, enkaz, hasar, y\i k\i m,\\
depremzede, 6 \c{S}ubat,\\
kahramanmara\c{s} depremi
\end{tabular} \\
\hline
Institutional / operational &
\begin{tabular}[t]{@{}l@{}}
AFAD, kurtarma, tahliye,\\
acil durum, \c{c}ad\i r kent,\\
bar\i nak, konteyner, prefabrik
\end{tabular} \\
\hline
Reconstruction &
\begin{tabular}[t]{@{}l@{}}
kal\i c\i\ konut, ge\c{c}ici bar\i nma,\\
yeniden yap\i m, yeniden yap\i lanma,\\
d\"on\"u\c{s}\"um, kentsel d\"on\"u\c{s}\"um,\\
TOK\.{I}, in\c{s}aat, altyap\i
\end{tabular} \\
\hline
\end{tabular}
\end{table}

This keyword pass should be understood as a deliberate scoping rule rather than a
neutral recovery mechanism for all earthquake-relevant discourse. It produces a
conservative analytic corpus centred on explicit disaster and reconstruction
vocabulary, but it may systematically exclude texts that discuss the earthquake
indirectly, metaphorically, or through less canonical local terminology. The
downstream results are therefore claims about this explicitly earthquake-marked
official discourse, not about every possible official text whose meaning may be
connected to the disaster.

% -----------------------------------------------------------------------------
\section{Annotation Codebook (Summary)}
\label{app:codebook}
% -----------------------------------------------------------------------------
The full codebook is included in the replication repository
(\texttt{annotation/codebook.md}). The operational core is summarised below.
Every frame is scored independently on a 0--3 ordinal scale; a single document may
score high on more than one frame.

\noindent\textbf{Decision procedure (applied per frame).}
\begin{enumerate}
\item \emph{Are the frame's indicators present at all?} No $\rightarrow$ 0.
\item \emph{Is the frame central or peripheral?} Peripheral (1--2 sentences)
      $\rightarrow$ 1.
\item \emph{Is the frame the backbone of the text?} Covers $\sim$30--50\%
      $\rightarrow$ 2; title + lead + conclusion all carry the frame $\rightarrow$ 3.
\end{enumerate}

\noindent\textbf{Frame definitions and indicators.}
\begin{itemize}
\item \textbf{Technical resilience} --- framing reconstruction through
engineering, building codes, inspection, materials, soil surveys, and damage
assessment. Indicators: \textit{yap\i\ denetimi, zemin et\"ud\"u, y\"onetmelik,
kentsel d\"on\"u\c{s}\"um} (engineering sense), \textit{g\"u\c{c}lendirme, AFAD
raporu}, numerical/technical detail.
\item \textbf{Political mobilisation} --- framing reconstruction through electoral
gain, partisan competition, leadership legitimacy, and national-mobilisation
rhetoric. Indicators: \textit{milletimize s\"oz verdik, muhalefet engelledi, Say\i n
Cumhurba\c{s}kan\i m\i z, AK Parti}, campaign pledges, ceremonial register.
\item \textbf{Developmental transformation} --- framing the disaster as a
development opportunity, investment, infrastructure modernisation, or regional
growth. Indicators: \textit{yeni \c{s}ehir, modern altyap\i, yat\i r\i m, TOK\.{I}
konutlar\i, OSB, istihdam}, project emphasis.
\item \textbf{Sustainability} --- framing reconstruction through climate
resilience, green building, green infrastructure, or ecological-city concepts;
rare in this corpus. Indicators: \textit{ye\c{s}il bina, iklim diren\c{c}li,
ekolojik, karbon ayak izi, yenilenebilir enerji, enerji verimlili\u{g}i}. Note:
``modern, comfortable housing'' is \emph{not} sustainability absent an explicit
environmental term.
\end{itemize}

% -----------------------------------------------------------------------------
\section{LLM Annotation Prompt}
\label{app:llmprompt}
% -----------------------------------------------------------------------------
The 201 documents outside the human double-coded subset were annotated via the
Anthropic API at temperature 0, with \texttt{max\_tokens} bounded and the model
asked to return a strict JSON object with the four frame scores. The prompt
combines (i) the codebook block reproduced verbatim below, (ii) eight few-shot
examples drawn from the human-agreement subset, and (iii) the target document
(truncated to a fixed character budget). The exact model identifier and API
version used are recorded in the configuration cell of the annotation notebook
(\texttt{annotation/LLM\_annotation.ipynb}). Two Anthropic models appear in the
notebook by design: an initial Claude Haiku 4.5 trial used for model selection,
and Claude Sonnet 4.5, which produced all 201 deployment-subset labels in the
final dataset. Haiku contributed no labels to the training set
(see Section~IV and the model-selection paragraph). The configuration cell pins
the exact model strings (\texttt{claude-sonnet-4-5} and
\texttt{claude-haiku-4-5-20251001}); these should match the prose verbatim.}

\noindent\textbf{Codebook block (verbatim, Turkish).}
\begin{quote}\footnotesize\ttfamily
4 frame'i 0-3 Likert \"ol\c{c}e\u{g}inde de\u{g}erlendir. Her frame BA\u{G}IMSIZ
de\u{g}erlendirilir; bir metin birden fazla frame'den y\"uksek puan alabilir.

FRAME 1 --- TECHNICAL; FRAME 2 --- POLITICAL; FRAME 3 --- DEVELOPMENT;
FRAME 4 --- SUSTAINABILITY (definitions and indicators as in
Appendix~\ref{app:codebook}).

KARAR REHBER\.{I}:
1. Frame'in g\"ostergeleri VAR MI? Hay\i r $\rightarrow$ 0. Evet:
2. Metnin MERKEZ\.{I}NDE M\.{I}? K\i y\i\ (1-2 c\"umle) $\rightarrow$ 1. Merkez:
3. Metnin OMURGASI MI? Hay\i r ($\sim$\%30-50) $\rightarrow$ 2.
   Evet (ba\c{s}l\i k+giri\c{s}+sonu\c{c}) $\rightarrow$ 3.
\end{quote}

\noindent The model was instructed to output only a JSON object of the form
\texttt{\{"technical": <0-3>, "political": <0-3>, "development": <0-3>,
"sustainability": <0-3>\}}, which was parsed and range-checked
(0$\le v \le$3) before being written to the label store with a
\texttt{label\_source} of \texttt{llm\_only} or, after manual review,
\texttt{llm\_then\_manual}.
 before \end{document} in discoure.tex
% Provides the replication appendices the reviewers asked for:
%   A. 26-term thematic keyword list (W1 / reproducibility)
%   B. Annotation codebook summary (reproducibility)
%   C. LLM annotation prompt + few-shot setup (reproducibility)
% =============================================================================

\appendices

% -----------------------------------------------------------------------------
\section{Thematic Keyword List}
\label{app:keywords}
% -----------------------------------------------------------------------------
The thematic keyword pass (Section~III.B) matches the following 26 terms
case-insensitively against the full cleaned text. A document is retained if it
matches at least one term. The list is reproduced in full here so that the
inclusion/exclusion decision is auditable; the prose summary in the body groups
these terms into three categories.

\begin{table}[!ht]
\renewcommand{\arraystretch}{1.2}
\caption{The 26 earthquake-relevance keywords used at the corpus-construction
stage. Matching is case-insensitive substring matching on the cleaned text.}
\label{table_keywords}
\centering
\footnotesize
\begin{tabular}{|l|l|}
\hline
\textbf{Category} & \textbf{Terms} \\
\hline
Core disaster &
\begin{tabular}[t]{@{}l@{}}
deprem, enkaz, hasar, y\i k\i m,\\
depremzede, 6 \c{S}ubat,\\
kahramanmara\c{s} depremi
\end{tabular} \\
\hline
Institutional / operational &
\begin{tabular}[t]{@{}l@{}}
AFAD, kurtarma, tahliye,\\
acil durum, \c{c}ad\i r kent,\\
bar\i nak, konteyner, prefabrik
\end{tabular} \\
\hline
Reconstruction &
\begin{tabular}[t]{@{}l@{}}
kal\i c\i\ konut, ge\c{c}ici bar\i nma,\\
yeniden yap\i m, yeniden yap\i lanma,\\
d\"on\"u\c{s}\"um, kentsel d\"on\"u\c{s}\"um,\\
TOK\.{I}, in\c{s}aat, altyap\i
\end{tabular} \\
\hline
\end{tabular}
\end{table}

This keyword pass should be understood as a deliberate scoping rule rather than a
neutral recovery mechanism for all earthquake-relevant discourse. It produces a
conservative analytic corpus centred on explicit disaster and reconstruction
vocabulary, but it may systematically exclude texts that discuss the earthquake
indirectly, metaphorically, or through less canonical local terminology. The
downstream results are therefore claims about this explicitly earthquake-marked
official discourse, not about every possible official text whose meaning may be
connected to the disaster.

% -----------------------------------------------------------------------------
\section{Annotation Codebook (Summary)}
\label{app:codebook}
% -----------------------------------------------------------------------------
The full codebook is included in the replication repository
(\texttt{annotation/codebook.md}). The operational core is summarised below.
Every frame is scored independently on a 0--3 ordinal scale; a single document may
score high on more than one frame.

\noindent\textbf{Decision procedure (applied per frame).}
\begin{enumerate}
\item \emph{Are the frame's indicators present at all?} No $\rightarrow$ 0.
\item \emph{Is the frame central or peripheral?} Peripheral (1--2 sentences)
      $\rightarrow$ 1.
\item \emph{Is the frame the backbone of the text?} Covers $\sim$30--50\%
      $\rightarrow$ 2; title + lead + conclusion all carry the frame $\rightarrow$ 3.
\end{enumerate}

\noindent\textbf{Frame definitions and indicators.}
\begin{itemize}
\item \textbf{Technical resilience} --- framing reconstruction through
engineering, building codes, inspection, materials, soil surveys, and damage
assessment. Indicators: \textit{yap\i\ denetimi, zemin et\"ud\"u, y\"onetmelik,
kentsel d\"on\"u\c{s}\"um} (engineering sense), \textit{g\"u\c{c}lendirme, AFAD
raporu}, numerical/technical detail.
\item \textbf{Political mobilisation} --- framing reconstruction through electoral
gain, partisan competition, leadership legitimacy, and national-mobilisation
rhetoric. Indicators: \textit{milletimize s\"oz verdik, muhalefet engelledi, Say\i n
Cumhurba\c{s}kan\i m\i z, AK Parti}, campaign pledges, ceremonial register.
\item \textbf{Developmental transformation} --- framing the disaster as a
development opportunity, investment, infrastructure modernisation, or regional
growth. Indicators: \textit{yeni \c{s}ehir, modern altyap\i, yat\i r\i m, TOK\.{I}
konutlar\i, OSB, istihdam}, project emphasis.
\item \textbf{Sustainability} --- framing reconstruction through climate
resilience, green building, green infrastructure, or ecological-city concepts;
rare in this corpus. Indicators: \textit{ye\c{s}il bina, iklim diren\c{c}li,
ekolojik, karbon ayak izi, yenilenebilir enerji, enerji verimlili\u{g}i}. Note:
``modern, comfortable housing'' is \emph{not} sustainability absent an explicit
environmental term.
\end{itemize}

% -----------------------------------------------------------------------------
\section{LLM Annotation Prompt}
\label{app:llmprompt}
% -----------------------------------------------------------------------------
The 201 documents outside the human double-coded subset were annotated via the
Anthropic API at temperature 0, with \texttt{max\_tokens} bounded and the model
asked to return a strict JSON object with the four frame scores. The prompt
combines (i) the codebook block reproduced verbatim below, (ii) eight few-shot
examples drawn from the human-agreement subset, and (iii) the target document
(truncated to a fixed character budget). The exact model identifier and API
version used are recorded in the configuration cell of the annotation notebook
(\texttt{annotation/LLM\_annotation.ipynb}). Two Anthropic models appear in the
notebook by design: an initial Claude Haiku 4.5 trial used for model selection,
and Claude Sonnet 4.5, which produced all 201 deployment-subset labels in the
final dataset. Haiku contributed no labels to the training set
(see Section~IV and the model-selection paragraph). The configuration cell pins
the exact model strings (\texttt{claude-sonnet-4-5} and
\texttt{claude-haiku-4-5-20251001}); these should match the prose verbatim.}

\noindent\textbf{Codebook block (verbatim, Turkish).}
\begin{quote}\footnotesize\ttfamily
4 frame'i 0-3 Likert \"ol\c{c}e\u{g}inde de\u{g}erlendir. Her frame BA\u{G}IMSIZ
de\u{g}erlendirilir; bir metin birden fazla frame'den y\"uksek puan alabilir.

FRAME 1 --- TECHNICAL; FRAME 2 --- POLITICAL; FRAME 3 --- DEVELOPMENT;
FRAME 4 --- SUSTAINABILITY (definitions and indicators as in
Appendix~\ref{app:codebook}).

KARAR REHBER\.{I}:
1. Frame'in g\"ostergeleri VAR MI? Hay\i r $\rightarrow$ 0. Evet:
2. Metnin MERKEZ\.{I}NDE M\.{I}? K\i y\i\ (1-2 c\"umle) $\rightarrow$ 1. Merkez:
3. Metnin OMURGASI MI? Hay\i r ($\sim$\%30-50) $\rightarrow$ 2.
   Evet (ba\c{s}l\i k+giri\c{s}+sonu\c{c}) $\rightarrow$ 3.
\end{quote}

\noindent The model was instructed to output only a JSON object of the form
\texttt{\{"technical": <0-3>, "political": <0-3>, "development": <0-3>,
"sustainability": <0-3>\}}, which was parsed and range-checked
(0$\le v \le$3) before being written to the label store with a
\texttt{label\_source} of \texttt{llm\_only} or, after manual review,
\texttt{llm\_then\_manual}.
